% De Lege Agraria III (de-lege-agraria-3) --- English translation
% Translated by Claude (Anthropic) from the Latin (Perseus).
% Style guide: STYLE.md.
%
% The speech survives only in part. The latter portion is lost.

\ciceroSection{1} The tribunes of the plebs would have done more becomingly, citizens, if what they bring before you about me they had said in my face rather, with me present; for they would have kept the fairness of your hearing, the practice of their predecessors, and the right of their own power. But since up to now they have shrunk from a present contest and contention, now, if it seems good to them, let them come forward in my assembly --- and the place to which they were unwilling to come when challenged by me, let them at least, called back, return to.

\ciceroSection{2} I see, citizens, certain men signifying I-know-not-what by their noise and presenting to my present assembly faces other than those they wore at my last. Therefore from you who have believed nothing about me I ask that you keep that goodwill toward me which you always had; from you whom I see somewhat altered, I ask the slight lease of a brief time of good opinion of me, so that, if I prove to you what I shall say, you may keep it forever; if otherwise, you may leave it deposited and cast away in this very place.

\ciceroSection{3} Your minds and ears, citizens, have been filled with the charge that I, doing a kindness to the Septimii, the Turranii, and the rest of the holders of Sullan assignments, am opposing the agrarian law and your interests. If anyone has believed this, he must first have believed this --- that by this agrarian law which has been promulgated the Sullan lands are taken away and divided among you, or at any rate that the holdings of private men are diminished so that you may be settled upon them. If I show that not only is no clod taken from anyone of the Sullan lands, but that this kind of land is most shamelessly confirmed and ratified by a particular clause of the law; if I prove that to those lands which were given by Sulla, Rullus by his law looks out so diligently that it readily appears the law was drafted not by a patron of your interests but by Valgius's son-in-law --- is there any reason, citizens, why he should not be seen by that charge which he used against me in my absence to have despised not my diligence and prudence only, but yours as well?

\ciceroSection{4} The fortieth clause of the law it is that I have, citizens, of set purpose not made mention of before you, lest I seem either to be rubbing again a scar already drawn over on the commonwealth, or to be stirring up at a most inopportune time some new dissension. Nor indeed shall I now dispute it because I think this present condition of the commonwealth not strongly to be defended --- particularly since I have professed myself to the Roman people for this year as the patron of quiet and concord --- but to teach Rullus to be silent hereafter at least about those matters about which he wishes silence to be kept concerning himself and his own deeds.

\ciceroSection{5} Of all laws I judge the most unjust and the most unlike a law to be the one which Lucius Flaccus the interrex carried about Sulla, that everything whatever Sulla had done should stand ratified. For when in other states, with tyrants set up, all laws are extinguished and abolished, here a tyrant of the commonwealth is set up by law. The law is invidious, as I said; yet it has its excuse: for it does not seem the law of a man, but of a time.

\ciceroSection{6} What if this is much more shameless? For by the Lex Valeria and by the Cornelian laws what is taken from a citizen is given to a citizen; an impudent favouring is joined with bitter injury; yet the man from whom it has been taken sips some hope from those laws, the man to whom it has been given some scruple. Rullus's provision is this: ``Those who after the consulship of Gaius Marius and Gnaeus Papirius . . .'' How far he flees from suspicion, in that he named precisely those consuls who were Sulla's chief adversaries! For if he had named Sulla as dictator, he thought it would be plain and invidious. But which of you did he think would be of so slow a wit that, after those consuls, the dictatorship of Sulla could not come to mind?

\ciceroSection{7} What then does the Marian tribune of the plebs say, who hales us Sullans into odium? ``Whoever after Marius and Carbo as consuls'' --- of land, buildings, lakes, ponds, places, possessions ---  he passes over the sky and the sea, embraces all the rest --- ``has had publicly given, assigned, sold, conceded'' --- by whom, Rullus? After Marius and Carbo as consuls, who assigned, who gave, who conceded, except Sulla? --- ``let all those things be of the same right'' --- of what right? he is, I suppose, undoing I-know-not-what; too sharp, too vehement a tribune of the plebs is annulling Sullan acts! --- ``as those things which are private under the best right.'' Even better than ancestral and inherited?

\ciceroSection{8} Better. But this the Lex Valeria does not say, the Cornelian laws do not ratify, Sulla himself does not demand. If those lands attain some part of right, some likeness of proper holding, some hope of long duration, no one of those holders is so shameless as not to think himself splendidly dealt with. But you, Rullus, what do you seek? that they should keep what they have? Who forbids it? that it should be private? It has so been carried. That your father-in-law's Hirpinian estate --- or the Hirpinian land, for he holds the whole of it --- should be of better right than my paternal and ancestral Arpinate estate?

\ciceroSection{9} For that is what you provide for. ``The best right'' belongs to those estates that are on the best terms. Free estates are of better right than burdened; by this clause everything which had been burdened will not be burdened. Discharged estates are in a better case than mortgaged; by the same clause everything mortgaged is, provided it is Sullan, freed. Untaxed are on a more advantageous condition than those which pay; I shall pay revenue to the Tusculans for the Crabra water, because I received the estate by mancipation; if it had been given me by Sulla, by Rullus's law I should not pay.

\ciceroSection{10} I see, citizens, that, as the matter itself compels you, you are stirred by the impudence either of the law or of the speech: of the law, which sets up a better right for Sullan estates than for ancestral; of the speech, which in such a case dares to charge anyone with too vehemently defending Sulla's accounts. Yet if it sanctioned only what had been given by Sulla, I should be silent, provided that he himself confessed himself a Sullan. But not only does he look out for those things, but he introduces another kind of donation entirely; and the man who charges me with defending Sullan possessions not only ratifies them but himself sets up new assignments and arises before us as a sudden Sulla.

\ciceroSection{11} Mark how great concessions of land this scolder of ours tries to make in a single word: ``which has been given, granted, conceded, sold.'' I let it pass, I hear it. What next? ``possessed.'' This a tribune of the plebs has dared to promulgate: that whatever a man possesses since the consulship of Marius and Carbo he should hold by that right by which what is private is held under the best title? Even if he expelled by force, even if he came into possession by stealth or by precarious tenure? By this law, then, the civil law, the suits about possessions, the praetors' interdicts, will be abolished.

\ciceroSection{12} No middling matter, citizens, no small theft lies hidden under this word. For there are many lands made public by the Cornelian law, neither assigned to anyone nor sold, that are most shamelessly held by a few men. For these he provides; these he defends; these he makes private. These lands, I say, which Sulla gave to no one, Rullus does not wish to assign to you but to bestow on those who hold them. I ask the reason why we should suffer to be sold what our ancestors left us in Italy, Sicily, Africa, the two Spains, Macedonia, Asia, when by the same law you see what is yours bestowed on the holders.

\ciceroSection{13} You will already see the whole law: that, since it was drafted for the domination of a few, it is also most accommodatingly fitted to the accounts of the Sullan assignment. For this man's father-in-law is a very good man; nor do I dispute now about his goodness, but about the son-in-law's impudence. He wishes to keep what he has and does not dissemble that he is a Sullan. This man, that he himself may have what he does not have, wishes through you to ratify what is doubtful, and, while he reaches for more than Sulla himself, charges me with defending Sullan accounts in the matters in which I resist.

\ciceroSection{14} ``He has,'' he says, ``some lands deserted and far away, my father-in-law has; he will sell them by my law for what price he wishes. He has some uncertain holdings, possessed by no right; they will be confirmed by the best right. He has public lands; I shall make them private. Finally, those estates which he has joined together as the best and most fruitful in the Casinate land --- proscribing the neighbours one after another so far as, by squaring the angles, he might form one boundary and norm of an estate out of many --- which he now holds with some fear, he will hold without any care.''

\ciceroSection{15} And since I have shown for what cause and for whose sake he has promulgated this, let him now teach in turn what holder I defend in resisting the agrarian law. You sell the Scantian Forest; the Roman people holds it; I defend. You divide the Campanian land; you are in possession; I do not yield. Then I see the holdings of Italy, Sicily, and the rest of the provinces put up for sale and proscribed by this law; they are your estates, your possessions; I shall resist and fight back, nor shall I allow the Roman people to be moved by anyone from his possessions in my consulship --- particularly, citizens, since nothing is sought for you.

\ciceroSection{16} For you ought no longer to be in this error. Are any of you fitted for violence, for crime, for slaughter? None. Yet for that kind of men, believe me, the Campanian land and that famous Capua are kept; an army is being set up against you, against your liberty, against Gnaeus Pompey; against this city Capua is set, against you the band of the most audacious men, against Gnaeus Pompey ten leaders are prepared. Let them come and face me, since they have called me out into your assembly at your demand, and let them dispute. \textit{[The remainder of the speech is lost.]}
